% =====================================================================
%  From Threats to Controls
%  A CIS-Grounded Solutions Dimension and Gameplay Framework
%  for the Security Cards
%
%  Prepared for: 30th Colloquium for Information Systems Security
%                Education (CISSE), submission deadline 15 Sept 2026.
%
%  IMPORTANT SUBMISSION NOTES
%  1. CISSE requires ANONYMIZED submissions. Author block below is
%     already blinded. Before submitting, also scrub: institution
%     names, program names, clinic names, IRB protocol numbers, and
%     the artifact URL (use an anonymized Zenodo/OSF view-only link).
%  2. Student authors must disclose student status at submission.
%  3. Use of generative AI tools must be disclosed; see the
%     disclosure block near the end of this file.
%
%  BUILD
%    pdflatex main
%    bibtex   main        (only if you switch to the BibTeX path below)
%    pdflatex main
%    pdflatex main
%
%  OVERLEAF / IEEEtran
%    To switch to the IEEE two-column style, comment out the
%    \documentclass line below and uncomment the IEEEtran line.
%    Everything else in this file is class-agnostic.
% =====================================================================

\documentclass[10pt,twocolumn,letterpaper]{article}
% \documentclass[conference]{IEEEtran}   % <-- Overleaf alternative

\usepackage[letterpaper,margin=0.75in,columnsep=0.25in]{geometry}
\usepackage{mathptmx}          % Times text + math
\usepackage[T1]{fontenc}
\usepackage[utf8]{inputenc}
\usepackage{graphicx}
\usepackage{booktabs}
\usepackage{tabularx}
\usepackage{array}
\usepackage{multirow}
\usepackage{xcolor}
\usepackage{enumitem}
\usepackage{caption}
\usepackage{float}
\usepackage{balance}
\usepackage{titlesec}
\usepackage[hidelinks,breaklinks]{hyperref}
\usepackage{url}

\captionsetup{font=small,labelfont=bf,skip=4pt}
\setlist{nosep,leftmargin=1.2em}
\titleformat{\section}{\normalfont\large\bfseries}{\thesection.}{0.5em}{}
\titleformat{\subsection}{\normalfont\normalsize\bfseries}{\thesubsection.}{0.5em}{}
\renewcommand{\arraystretch}{1.15}
\newcommand{\ig}[1]{\textsc{ig#1}}
\newcommand{\todo}[1]{\textcolor{red}{[\textbf{TODO:} #1]}}

% Placeholder box for figures whose image files are not yet dropped in.
\newcommand{\figplaceholder}[2]{%
  \fbox{\parbox[c][#1][c]{0.96\linewidth}{\centering\small\itshape #2}}%
}

\title{\vspace{-1.5em}\bfseries From Threats to Controls: A CIS-Grounded
Solutions Dimension and Gameplay Framework for the Security Cards}

\author{\normalsize Anonymous Submission\\
\normalsize \textit{30th Colloquium for Information Systems Security Education}}
\date{}

\begin{document}
\maketitle

% =====================================================================
\begin{abstract}
\noindent\small
The Security Cards, developed at the University of Washington in 2013, are a
widely used toolkit for broadening how learners think about security threats
across four dimensions: Human Impact, Adversary's Motivations, Adversary's
Resources, and Adversary's Methods. The toolkit deliberately stops at threat
discovery. In cybersecurity clinic courses, where students advise
under-resourced community organizations, discussion that ends at threat
identification leaves the central professional task---recommending
proportionate controls under real budget constraints---unpracticed. We present
a fifth dimension, \emph{Solutions}, comprising twenty-one defensive control
cards mapped to the CIS Critical Security Controls v8.1 at the safeguard level,
together with seven zero-cost modifier cards that carry architectural and
human-factors principles rather than controls. Card cost is
derived from CIS Implementation Group tier rather than assigned by designer
judgment, which makes budgeted play a defensible model of resource-constrained
decision making. A coverage analysis against \ig{1} identified five foundational
gaps in our initial fifteen-card deck---asset inventory, software inventory,
secure configuration, email and browser protections, and incident response
planning---which motivated the current version. We contribute the mapped deck,
a six-variant gameplay framework spanning increasing cognitive demand, and a
three-phase core loop in which a risk matrix is scored before and after control
selection so that residual risk becomes both a game mechanic and a measurable
artifact. A formative evaluation with 19 multidisciplinary participants
established feasibility and accessibility; we make no claim of learning gains.
A controlled classroom study is scheduled for September 2026 and its design is
reported in full. All cards are released under an open license.

\vspace{0.6em}
\noindent\textbf{Keywords:} cybersecurity education, game-based learning,
serious games, CIS Controls, risk assessment, cybersecurity clinics,
security controls, tabletop exercise
\end{abstract}

\vspace{0.8em}

% =====================================================================
\section{Introduction}

Cybersecurity education faces a persistent structural problem: the material
changes faster than curricula can be rewritten, and the reasoning that
practitioners actually perform---weighing partial information, arguing for a
position, and committing resources under constraint---is poorly served by
lecture and by multiple-choice assessment. Active and collaborative methods
address this gap in engineering education generally \cite{prince2004active},
and game-based learning has emerged as one practical vehicle in security
specifically \cite{plass2015gbl,wouters2013meta}.

Our setting is the cybersecurity clinic. Clinics are skill-based programs in
which university students deliver cybersecurity services to under-resourced
organizations in their community \cite{clinics_consortium}. Clinic students face
an unusual demand curve. Within a single semester they move from little or no
security background to sitting across a table from a nonprofit director, a
municipal clerk, or a small clinic administrator and recommending what that
organization should do next, given a budget that is frequently near zero. The
early weeks of the curriculum therefore need an activity that (i) requires no
prerequisite technical vocabulary, (ii) exercises prioritization rather than
recall, and (iii) ends in a defensible recommendation rather than a list of
worries.

The Security Cards \cite{denning2013securitycards} solve the first two
requirements well. Developed by the University of Washington's Security and
Privacy Research Lab and Value Sensitive Design Research Lab, the 42-card
toolkit organizes security thinking into four dimensions and is explicitly
designed to widen the space of threats a group considers. It is, by its
authors' framing, a brainstorming toolkit rather than a game: there is no win
condition, no scoring, and no resource constraint. This is a deliberate and
appropriate design choice for its stated purpose.

It is also the source of the third gap. In our own use, card-driven discussion
reliably produced rich threat enumeration and then stopped. Students left the
table with a longer list of things that could go wrong and no practiced
vocabulary for what to do about it. The original authors anticipated this and
suggested a ``Considering Defense'' activity, but supplied no defense cards, no
control framework, and no mechanism for making defensive choices costly.

This paper contributes the missing half of that loop:

\begin{enumerate}[label=(\arabic*)]
\item \textbf{A Solutions dimension of twenty-eight cards}---twenty-one
defensive controls and six principle-based modifiers---mapped to CIS
Critical Security Controls v8.1 \cite{cis2024controls} at the level of
individual safeguards, not merely control titles, so that a novice can see what
``access control management'' concretely requires.
\item \textbf{A cost model derived from CIS Implementation Groups}
\cite{cis_ig} rather than from designer intuition. \ig{1} controls cost one
point, \ig{2} two, \ig{3} three. Budgeted play then reproduces the actual
economics of the small organizations clinic students serve.
\item \textbf{A coverage analysis} of our deck against \ig{1}, which
identified five foundational omissions and supplied a non-arbitrary rationale
for the current card set. We report the analysis rather than the intuition
that preceded it.
\item \textbf{A six-variant gameplay framework} of increasing cognitive
demand, and a three-phase core loop---threat construction, risk placement,
budgeted defense, residual-risk replacement---in which the risk matrix is
scored twice and the movement between placements is simultaneously the
pedagogical point and a measurable outcome.
\item \textbf{A formative evaluation} with 19 multidisciplinary participants
establishing feasibility, accessibility, and time-to-play, together with the
full design of a controlled classroom study scheduled for September 2026.
\end{enumerate}

We are explicit about what this paper does not claim. The evaluation reported
in Section~\ref{sec:eval} is formative. It establishes that the artifact is
usable by novices in under ten minutes and that participants found it
approachable. It does not establish learning gains, and we do not argue that it
does. Section~\ref{sec:planned} states the study that would.

Section~\ref{sec:background} situates the work among existing security card and
tabletop games. Section~\ref{sec:design} presents the deck, the safeguard
mapping, and the coverage analysis. Section~\ref{sec:gameplay} presents the
gameplay framework. Section~\ref{sec:walkthrough} walks through two recorded
sessions. Sections~\ref{sec:eval} through \ref{sec:discussion} report the
formative evaluation, the planned study, and limitations.

% =====================================================================
\section{Background and Related Work}
\label{sec:background}

\subsection{Terminology}

The literature distinguishes two ideas that are frequently conflated, including
in our own earlier drafts. \emph{Gamification} is the use of game design
elements---points, badges, leaderboards---in contexts that are not themselves
games \cite{deterding2011gamification}. \emph{Game-based learning} (GBL), by
contrast, uses an actual game as the learning environment
\cite{plass2015gbl}; Landers' theory of gamified learning
\cite{landers2014gamified} clarifies that the two operate through different
mechanisms, gamification by moderating behavior toward the instructional
content and GBL by making the content itself the object of play.

This paper is GBL work. We add rules, roles, scarcity, and scoring to a toolkit
that previously had none, which is why the distinction matters for our
contribution rather than being a matter of vocabulary hygiene: the claim is
that turning a brainstorming aid into a game with constraints changes what
learners do at the table, not that adding points to a lecture increases
attention.

\subsection{Card and Tabletop Games in Security Education}

Security education has a well-developed tradition of physical card and tabletop
games, catalogued extensively by Shostack \cite{shostack_games}. Several are
directly relevant.

\textbf{Control-Alt-Hack} \cite{denning2013controlalthack} targets a novice and
general audience, using a fictional penetration-testing consultancy to build
awareness of what security work involves. Its aim is dispositional---conveying
that the field exists and is interesting---rather than teaching a specific
analytic procedure.

\textbf{Elevation of Privilege} \cite{shostack2014eop} converts STRIDE threat
modeling into a trick-taking card game and demonstrably draws non-specialists
into structured threat analysis. It assumes a system diagram and a developer
audience, which places it downstream of where clinic students begin.

\textbf{[d0x3d!]} \cite{gondree2013d0x3d} is a cooperative board game about
network assets and adversary movement, aimed at building intuition for
attacker--defender dynamics.

\textbf{Riskio} \cite{hart2020riskio} is a tabletop game for non-technical
staff in organizations, in which players defend organizational assets against
STRIDE-categorized threats. It requires an experienced facilitator acting as
Game Master.

\textbf{Backdoors \& Breaches} \cite{bhis2019backdoors} is the closest analogue
to our Solutions dimension: players draw hidden attack cards and select from a
fixed set of procedure cards to detect and respond. It has been used
successfully in graduate and undergraduate classrooms. Its focus is incident
response---detecting an attack already underway---rather than \emph{ex ante}
control selection, and its procedure cards are detection techniques rather than
a framework-aligned control catalogue.

\textbf{MalAware} \cite{malaware2024} structures a malware tabletop exercise
around the NIST Cybersecurity Framework functions, again with an incident
response emphasis.

\textbf{Cyber Defence Dice} \cite{cyberdefencedice2026}, published recently in
this journal, pursues a complementary goal: reducing setup and play time to the
minimum required for entry-level awareness raising, and its survey of the field
identifies facilitator dependence and session length as recurring barriers.

Table~\ref{tab:related} positions these works along the two axes that matter
for our setting: whether the game addresses threat identification, control
selection, or incident response; and whether it is usable by learners with no
prior security background and no expert facilitator.

\begin{table*}[t]
\centering
\caption{Positioning of the proposed deck among existing security card and tabletop games.}
\label{tab:related}
\small
\begin{tabularx}{\textwidth}{@{}l l l l X@{}}
\toprule
\textbf{Artifact} & \textbf{Primary focus} & \textbf{Audience} &
\textbf{Framework anchor} & \textbf{Relation to this work} \\
\midrule
Security Cards \cite{denning2013securitycards} & Threat discovery & Novice--expert & None &
Base artifact; four dimensions retained unchanged \\
Control-Alt-Hack \cite{denning2013controlalthack} & Awareness / disposition & Novice & None &
Different goal; no analytic procedure taught \\
Elevation of Privilege \cite{shostack2014eop} & Threat modeling & Developers & STRIDE &
Assumes a system diagram and technical vocabulary \\
{[d0x3d!]} \cite{gondree2013d0x3d} & Attacker--defender dynamics & Novice & None &
Board-based; abstracted network model \\
Riskio \cite{hart2020riskio} & Threat and asset defense & Non-technical staff & STRIDE &
Requires an expert Game Master \\
Backdoors \& Breaches \cite{bhis2019backdoors} & Incident response & Mixed & Kill chain &
Closest analogue; response rather than \emph{ex ante} control selection \\
MalAware \cite{malaware2024} & Malware incident response & Mixed & NIST CSF &
Response emphasis; single threat family \\
Cyber Defence Dice \cite{cyberdefencedice2026} & Rapid awareness raising & Novice & --- &
Complementary; optimizes for minimal play time \\
\midrule
\textbf{This work} & \textbf{Threat $\rightarrow$ control mapping} & \textbf{Novice} &
\textbf{CIS Controls v8.1 (\ig{1})} &
\textbf{Budgeted control selection with residual-risk scoring} \\
\bottomrule
\end{tabularx}
\end{table*}

The gap this table makes visible is specific. Existing games teach learners to
\emph{find} threats or to \emph{respond} to incidents. None asks a novice to
select a proportionate set of preventive controls against a stated budget and
justify the choice against a recognized framework. That is precisely the task a
clinic student performs for a client in week eight.

\subsection{The CIS Controls and Implementation Groups}

The CIS Critical Security Controls v8.1 organize defensive practice into 18
controls containing 153 individual safeguards \cite{cis2024controls}. Crucially
for our purpose, CIS partitions the safeguards into three Implementation Groups
\cite{cis_ig}. \ig{1} defines essential cyber hygiene: the 56 safeguards that a
small organization with limited expertise should implement first. \ig{2} adds
safeguards for organizations with dedicated security staff. \ig{3} addresses
organizations with specialist capability.

Three properties make this framework an unusually good fit for clinic
education. First, \ig{1} is defined \emph{for} the class of organization clinic
students serve, so the deck teaches the right subset rather than an arbitrary
one. Second, the tiering supplies an externally justified difficulty and cost
ordering, removing the need for designer-assigned point values. Third, control
identifiers give students a durable reference: a student who learns that
``patch things'' is CIS Control 7 can find the safeguards, the measures, and
the vendor mappings after the game ends.

% =====================================================================
\section{Design of the Solutions Dimension}
\label{sec:design}

\subsection{Design Goals}

Four goals governed the design, the first two inherited from the original
toolkit and the last two specific to the clinic setting:

\begin{description}[leftmargin=1em,style=nextline]
\item[G1. Vocabulary accessibility.] A participant with no prior security
coursework should be able to read a card and act on it without external
explanation.
\item[G2. Discussion generation.] Cards should provoke disagreement, which is
the mechanism by which the activity teaches, rather than converge quickly on a
single obvious answer.
\item[G3. Framework traceability.] Every defensive claim made at the table
should be traceable to a published safeguard, so that instructors can extend
the discussion and students can continue independently.
\item[G4. Constraint realism.] Choosing a control should cost something, because
the defining feature of the organizations clinic students serve is that they
cannot do everything.
\end{description}

\subsection{Revisions to the Original Card Format}

We preserved the double-sided format, the dimension coloring scheme, and the
example-driven card backs of the original deck. Two changes were made to the
existing dimensions.

\textbf{Iconography.} The original cards carry an evocative photograph intended
to prompt lateral association. In our pilot sessions, facilitators observed
participants asking what the photograph depicted and interpreting it
literally---for example, reading the ventilation-grille image on the
\emph{Physical Attack} card as being about building HVAC systems specifically.
We replaced photographs with flat pictographic icons chosen to denote the card
topic rather than to evoke around it. We report this as a design decision
motivated by facilitator observation during unstructured pilots, not as an
empirical finding: we did not measure comprehension time or accuracy under the
two conditions, and the icon change is confounded with other simultaneous
revisions. Section~\ref{sec:planned} includes a card-comprehension check
designed to test it properly.

\textbf{An Artificial Intelligence resource card.} Generative models have
materially lowered the cost of producing convincing phishing content, cloned
voice and video, and target reconnaissance. Because the Adversary's Resources
dimension enumerates what an attacker can bring to bear, we added an
\emph{Artificial Intelligence} card to that dimension rather than treating AI
as a separate threat category. This keeps the original dimensional logic
intact.

\textbf{Color separation.} Our first Solutions cards used a blue palette that
sat close to the existing Human Impact blue, and pilot participants
occasionally sorted the two together. The current Solutions cards use a
distinct hue \emph{and} a distinct icon frame shape, so that the dimensions
remain separable in grayscale printing and for color-vision-deficient players.
\todo{Confirm final hex values and include a grayscale proof in the camera-ready.}

\subsection{Solutions Cards and Safeguard Mapping}

Each Solutions card names a defensive practice, states what it does in one
sentence of non-specialist language, gives two concrete examples, and cites the
CIS Control and representative safeguards it implements.
Table~\ref{tab:mapping} presents the complete mapping; Appendix~\ref{app:cards}
gives the full text of all twenty-eight cards, including the discussion prompt
each one carries.

Two properties of the table deserve comment. First, mapping is at the safeguard
level. A card that said only ``CIS Control 6: Access Control Management'' would
tell a novice nothing; naming safeguards 6.1 and 6.2---grant access on a
documented process, revoke it when a person leaves---tells them what the control
actually asks an organization to do. Second, the \ig{} column and the cost
column are the same information. Cost is not an aesthetic judgment; it is the
Implementation Group at which the card's representative safeguards sit. Where a
card spans tiers, cost reflects the tier of the implementation a small
organization would realistically adopt, and the card back states the tier range.

\begin{table*}[t]
\centering
\caption{Solutions dimension: cards, CIS Control v8.1 mapping, representative safeguards,
Implementation Group, and derived play cost. Cards 9--13 were added in the current
version following the \ig{1} coverage analysis in Section~\ref{sec:gaps}.}
\label{tab:mapping}
\scriptsize
\begin{tabularx}{\textwidth}{@{}r l l X c c@{}}
\toprule
\textbf{\#} & \textbf{Card} & \textbf{CIS Control} & \textbf{Representative safeguards} &
\textbf{IG} & \textbf{Cost} \\
\midrule
\multicolumn{6}{@{}l}{\textit{Foundational tier (cost 1)}}\\
\midrule
1  & Software Updates              & 7 Continuous Vulnerability Mgmt & 7.1 establish process; 7.3 automated OS patching; 7.4 automated application patching & 1 & 1 \\
2  & Social Engineering Awareness  & 14 Security Awareness \& Skills Training & 14.1 establish program; 14.2 recognize social engineering; 14.3 authentication best practice & 1 & 1 \\
3  & Deepfake \& AI Fraud Awareness & 14 Security Awareness \& Skills Training & 14.2 extended to synthetic media; 14.9 role-specific training (\ig{2}) & 1--2 & 1 \\
4  & Identity \& Access Management & 6 Access Control Management & 6.1 documented access-grant process; 6.2 access revocation; 6.7 centralized access control (\ig{2}) & 1--2 & 1 \\
5  & Password \& Authentication    & 5 Account Mgmt; 6 Access Control Mgmt & 5.2 unique passwords; 5.4 dedicated admin accounts; 6.3--6.5 multi-factor authentication & 1 & 1 \\
6  & Backup \& Recovery            & 11 Data Recovery & 11.1 recovery process; 11.2 automated backups; 11.3 protect recovery data; 11.4 isolated copy & 1 & 1 \\
7  & Account Lifecycle Management  & 5 Account Management & 5.1 account inventory; 5.3 disable dormant accounts; 5.5 service account inventory (\ig{2}) & 1--2 & 1 \\
8  & Physical Security             & --- (supports 1, 3) & No dedicated v8.1 control; enables asset and data safeguards & --- & 1 \\
9  & Asset Inventory               & 1 Inventory \& Control of Enterprise Assets & 1.1 maintain asset inventory; 1.2 address unauthorized assets & 1 & 1 \\
10 & Software Inventory            & 2 Inventory \& Control of Software Assets & 2.1 software inventory; 2.2 ensure software is supported; 2.3 address unauthorized software & 1 & 1 \\
11 & Secure Configuration          & 4 Secure Configuration & 4.1 configuration process; 4.2 network infrastructure config; 4.6 secure management; 4.7 manage default accounts & 1 & 1 \\
12 & Email \& Browser Protections  & 9 Email \& Web Browser Protections & 9.1 supported browsers and clients; 9.2 DNS filtering; 9.6 block unnecessary file types (\ig{2}) & 1--2 & 1 \\
13 & Incident Response Planning    & 17 Incident Response Management & 17.1 designate personnel; 17.2 reporting contacts; 17.3 reporting process; 17.4 handling process (\ig{2}) & 1--2 & 1 \\
\midrule
\multicolumn{6}{@{}l}{\textit{Extended tier (cost 2)}}\\
\midrule
14 & Encryption                    & 3 Data Protection & 3.6 encrypt end-user devices (\ig{1}); 3.10 encrypt data in transit; 3.11 encrypt data at rest & 1--2 & 2 \\
15 & Audit Logging \& SIEM         & 8 Audit Log Mgmt; 13 Network Monitoring & 8.2 collect audit logs (\ig{1}); 8.9 centralize logs; 8.11 conduct log reviews; 13.1 centralized alerting & 1--2 & 2 \\
16 & Endpoint Detection \& Response & 10 Malware Defenses; 13 Network Monitoring & 10.1 anti-malware (\ig{1}); 10.6 centrally managed; 10.7 behavior-based detection & 1--2 & 2 \\
17 & Third-Party \& Contractual Controls & 15 Service Provider Management & 15.1 service provider inventory (\ig{1}); 15.4 security requirements in contracts; 15.5 assessment (\ig{3}) & 1--3 & 2 \\
18 & Network Segmentation          & 12 Network Infrastructure Mgmt; 13 & 12.2 secure network architecture; 13.4 traffic filtering between segments & 2 & 2 \\
\midrule
\multicolumn{6}{@{}l}{\textit{Advanced tier (cost 3)}}\\
\midrule
19 & Penetration Testing           & 18 Penetration Testing & 18.1 establish program; 18.2 external tests; 18.3 remediate findings; 18.5 internal tests (\ig{3}) & 2--3 & 3 \\
20 & Managed Detection \& Response & 8 Audit Log Mgmt; 13 Network Monitoring; 17 Incident Response & 8.11 conduct log reviews; 13.1 centralized alerting; 13.11 tune alerting thresholds (\ig{3}); 17.1 designate incident personnel & 2--3 & 3 \\
21 & Cyber Insurance \& Risk Transfer & --- (risk treatment, not a control) & No CIS control covers transfer; insurers increasingly require \ig{1} safeguards as a condition of coverage & --- & 3 \\
\midrule
\multicolumn{6}{@{}l}{\textit{Modifier cards (cost 0; three dealt per round, at most two played)}}\\
\midrule
22 & Defense in Depth              & Architectural principle & Requires two purchases that fail by different means & --- & 0 \\
23 & Fail Safe Defaults            & Architectural principle & Requires naming the state a purchased control fails into & --- & 0 \\
24 & Least Privilege               & Architectural principle & Requires card 4 or 7; name one access to reduce & --- & 0 \\
25 & Human Factors                 & Socio-technical principle & Requires naming a purchase staff will circumvent & --- & 0 \\
26 & Risk Acceptance               & Risk treatment (extended) & Played instead of a purchase; requires a named decision owner & --- & 0 \\
27 & Compensating Control          & Risk treatment (extended) & Requires naming the unaffordable control and the residual gap & --- & 0 \\
28 & Deferred Investment           & Multi-engagement planning (extended) & Caps spend at 4; banks 2 for a 10-point budget next engagement, with accrued risk & --- & 0 \\
\bottomrule
\end{tabularx}
\end{table*}

\subsection{Coverage Analysis and the Rationale for Version 2}
\label{sec:gaps}

An early reviewer of this work asked a question we could not answer at the
time: on what basis were these particular fifteen cards chosen? The honest
answer was designer intuition informed by clinic experience. To replace
intuition with a procedure, we performed a coverage analysis: for each of the
56 \ig{1} safeguards, we asked whether any card in the deck would prompt a
player to consider it.

The result was uncomfortable and useful. Our initial deck omitted five
foundational areas while including one that sits outside \ig{1} entirely:

\begin{itemize}
\item \textbf{Controls 1 and 2 (asset and software inventory).} The first two
controls in the framework, and arguably the precondition for every other one,
were absent. An organization cannot patch software it does not know it runs.
\item \textbf{Control 4 (secure configuration).} Absent, despite default
credentials and unhardened configurations being among the most common findings
in clinic engagements.
\item \textbf{Control 9 (email and web browser protections).} Absent, even
though phishing was the central threat in both of our recorded walkthrough
sessions. Students repeatedly reached for awareness training as the only
anti-phishing lever available in the deck because it was the only one present.
\item \textbf{Control 17 (incident response planning).} Absent, despite clinic
students being the people a client calls during an incident.
\item \textbf{Control 18 (penetration testing) was present} and is \ig{2}/\ig{3}
only---inconsistent with our stated \ig{1} orientation. We retained it, because
clinic students are asked about it by clients, but re-tiered it to cost~3 so
that its expense is visible in play.
\end{itemize}

Cards 9 through 13 in Table~\ref{tab:mapping} close these gaps. Every added card
is \ig{1} at cost~1, which shifts the deck's center of gravity toward essential
cyber hygiene, as originally intended but not originally achieved. We
deliberately did not add a card for Control 16 (Application Software Security),
which begins at \ig{2} and addresses in-house software development---an activity
the client organizations in question rarely undertake.

We report this analysis because the procedure, not the resulting card list, is
the transferable contribution. Any deck extension of this kind can be audited
against a published framework rather than defended on the designer's taste.

\subsection{Calibrating the Budget}
\label{sec:budget}

A budget constrains only if it is small relative to what is on the table. Our
first implementation dealt eight cards against ten points, which on inspection
is barely a constraint at all. The twenty-one purchasable cards carry 32 points
in total, a mean of 1.52 points per card, so a random eight-card hand is worth
approximately 12.2 points. Ten points buys 82\% of such a hand, and a hand drawn
entirely from the Foundational tier is fully purchasable. A team in that
position makes no trade-off; it simply buys everything and the mechanic teaches
nothing.

We therefore set the default budget at \textbf{six points}, which purchases
roughly half a typical hand and makes the single Advanced-tier card cost half of
everything a team has. Table~\ref{tab:budget} shows the calibration.

\begin{table}[h]
\centering
\caption{Budget calibration against an eight-card hand worth approximately 12.2
points. Six points is the default; larger budgets model better-resourced
organizations.}
\label{tab:budget}
\footnotesize
\begin{tabular}{@{}c c c l@{}}
\toprule
\textbf{Budget} & \textbf{\% hand} & \textbf{Cost-3 card} & \textbf{Organization modeled} \\
\midrule
5  & 41\% & 60\% & Volunteer-run, no IT staff \\
\textbf{6}  & \textbf{49\%} & \textbf{50\%} & \textbf{Default; \ig{1}} \\
10 & 82\% & 30\% & \ig{2}; one IT staff member \\
15 & --- & 20\% & \ig{3}; security function \\
\bottomrule
\end{tabular}
\\[0.3em]{\footnotesize ``Cost-3 card'' is the share of budget one Advanced-tier card consumes.}
\end{table}

The answer to whether the budget should ever exceed six is therefore yes, but
not as a general relaxation. Budget is a property of the \emph{organization} in
the scenario, not of the game, and tying it to Implementation Group makes that
explicit: the same threat facing a volunteer-run food bank and a county health
department warrants different answers, and students should feel that difference
rather than be told about it. At fifteen points the pedagogical problem changes
character---a team can afford nearly everything, so the interesting question
becomes sequencing rather than selection, which is a more advanced skill worth
reaching later in a curriculum.

For the controlled study in Section~\ref{sec:planned} we hold the budget fixed
at six across all teams, since varying it would confound the condition
comparison.

\subsection{Managing Deck Growth}

Growing from fifteen to twenty-eight cards runs against a finding from our own pilot,
in which several participants reported that fifteen Solutions cards were already
a lot to hold in mind. Two mechanisms address this without shrinking the
catalogue.

First, \emph{the deck is not the hand}. No gameplay variant deals more than
eight control cards and three modifiers at once. The catalogue exists so that instructors can
compose scenario-appropriate subsets; players never face the full catalogue.

Second, \emph{the tiers are a curriculum sequence, not a difficulty label}.
Foundational-tier cards alone (thirteen cards, all cost~1) support the first
weeks of a clinic course. The Extended and Advanced tiers are introduced later,
once students have a reason to ask why anyone would spend three points on a
single control. This mirrors the progression CIS itself describes for
organizations maturing from \ig{1} to \ig{2}.

\subsection{Modifier Cards and the Justification Requirement}
\label{sec:modifiers}

The seven modifier cards are a distinct card class. They carry principles rather
than controls, they cost nothing, and each is governed by the same rule: a
modifier may be played only if the team can articulate the specific
justification the card demands. \emph{Defense in Depth} requires naming two
purchases that fail by different means; \emph{Fail Safe Defaults} requires
naming the state a purchased control fails into \cite{saltzer1975protection};
\emph{Least Privilege}
requires naming one access the team would reduce and what would break. If the
team cannot supply the justification, the card returns to the deck.

This constraint is what keeps a zero-cost card from being a free point. It also
serves the evaluation directly: because every modifier requires a spoken and
recorded rationale, modifier play generates precisely the artifacts that measure
M3 in Section~\ref{sec:planned} scores. Three of the seven---\emph{Risk Acceptance}, \emph{Compensating Control}, and
\emph{Deferred Investment}---encode risk treatments and planning horizons other
than immediate mitigation. They are conceptually harder than the rest of the
deck, so we assign them to the Extended tier rather than to introductory play.

The class exists because the original deck, and our first version of the
Solutions dimension, taught exactly one risk treatment: buy a control. Practice
recognizes four---avoid, mitigate, transfer, accept---and a student who leaves a
clinic course believing that every risk must be fixed will give clients advice
they cannot afford to take. \emph{Cyber Insurance} supplies transfer,
\emph{Risk Acceptance} supplies acceptance, and \emph{Compensating Control}
supplies the partial mitigation that under-resourced organizations actually
implement.

\subsection{Availability}

The full deck is available as print-ready files under a Creative Commons
license, with a persistent identifier.
\todo{Mint a Zenodo DOI and replace the Canva link, which was reported broken by
a reviewer. For the anonymized submission, substitute an anonymous
view-only link.}

% =====================================================================
\section{Gameplay Framework}
\label{sec:gameplay}

The original toolkit describes activities but not rules. We specify six
variants with explicit objectives, durations, and produced artifacts
(Table~\ref{tab:variants}). The variants are ordered by cognitive demand and are
designed to be run in sequence across a semester or selected individually to fit
an available class period.

\begin{table*}[t]
\centering
\caption{Gameplay variants, ordered by cognitive demand. ``Artifact'' names the
observable product of play, which is what makes each variant assessable.}
\label{tab:variants}
\small
\begin{tabularx}{\textwidth}{@{}l X l l X@{}}
\toprule
\textbf{Variant} & \textbf{Objective} & \textbf{Cards used} & \textbf{Time} &
\textbf{Artifact produced} \\
\midrule
V1 Threat Model Generation & Construct a coherent incident narrative from
component threat elements & One card from each of the four original dimensions
& 8--10 min & Written incident narrative \\
V2 Risk Matrix Placement & Estimate likelihood and severity; reconcile
disagreement & Adversary's Methods (7 cards) & 10--15 min & Matrix coordinates
per threat, with rationale \\
V3 Solution Ranking & Prioritize among defensible alternatives & 8 dealt
Solutions cards & 7--10 min & Ranked top-three with justifications \\
V4 Control Mapping & Attribute an incident to the control that failed &
Solutions + a stated scenario & 10 min & Control attribution with reasoning \\
V5 Budgeted Defense & Allocate scarce resources across a control portfolio &
Solutions + 6-point budget & 15--20 min & Purchased portfolio and spend
rationale \\
V6 Residual Risk Loop & Recognize that controls reduce rather than eliminate
risk & Full sequence, V2 $\rightarrow$ V5 $\rightarrow$ V2 & 20--25 min &
Pre/post matrix placements and the delta between them \\
\bottomrule
\end{tabularx}
\end{table*}

\subsection{The Core Loop}

V6 is the flagship variant and the one we recommend for the classroom study.
It runs in three phases.

\textbf{Phase 1 --- Threat model and placement.} Each team draws one card from each of
the four original dimensions and constructs a scenario in the manner of V1. The
team then places the resulting threat on a $3 \times 3$ risk matrix
(Figure~\ref{fig:matrix}) and records the coordinates. Placement on the severity
axis must be justified by naming the \emph{Human Impact} card in play and saying
who specifically is harmed and how badly. This constraint matters more than it
appears: novices reliably rate severity by how technically sophisticated an
attack sounds rather than by its consequences, and requiring them to point at
the Human Impact card forces severity to mean what it means in practice.

\textbf{Phase 2 --- Budgeted defense.} The team receives eight Solutions cards
and a budget of six points. It purchases a portfolio and states, for each
purchase, which safeguard addresses which element of the threat. Three modifier
cards are dealt alongside the eight; at most two may be played. The
\emph{Defense in Depth} modifier card costs nothing but may only be played
alongside two other purchased cards addressing the same threat by different
means---an inexpensive mechanic that teaches layering without requiring a
lecture on it.

\textbf{Phase 3 --- Residual risk.} The team re-places the same threat on the
matrix, given the controls it has purchased, and states what would still get
through. Teams then compare their pre- and post-placements with those of other
teams.

This structure does three things at once. Pedagogically, the second placement is
where students discover that controls shift risk rather than remove it, and
where the group that spent all ten points on prevention has to explain how it
would detect the attack it did not prevent. Mechanically, it closes the loop
between risk assessment and control selection that a reviewer of our earlier
draft correctly identified as unexplained. Methodologically, the pre/post delta
and the accompanying justifications are recorded artifacts that can be scored
against an expert key, which converts a discussion activity into a measurable
one without adding a test.

\begin{figure}[t]
\centering
\small
\begin{tabular}{|l|c|c|c|}
\hline
 & \textbf{Not Severe} & \textbf{Severe} & \textbf{Very Severe} \\
\hline
\textbf{Not Likely}  & \rule{0pt}{2.2em} & & \\
\hline
\textbf{Likely}      & \rule{0pt}{2.2em} & & \\
\hline
\textbf{Very Likely} & \rule{0pt}{2.2em} & & \\
\hline
\end{tabular}
\caption{The $3 \times 3$ risk matrix used in variants V2 and V6. In V6 the same
threat is placed twice: once before and once after control selection.}
\label{fig:matrix}
\end{figure}

\subsection{Campaign Play Across a Semester}
\label{sec:campaign}

A single round teaches selection under constraint. It does not teach what
clinic students actually deliver, which is a multi-year roadmap: a statement of
what the organization should do this year, what it should defer, and what the
deferral costs while it waits.

The \emph{Deferred Investment} modifier (card 28) introduces this. A team may
cap its spend at four of six points and bank the remaining two; in the next
engagement it receives ten points rather than six. The multiplier is deliberate:
ten points is the \ig{2} budget in Table~\ref{tab:budget}, so underspending now
is precisely how an organization climbs from essential cyber hygiene to the next
Implementation Group. That is the CIS maturity model expressed as a game
mechanic rather than described in a lecture.

Deferral must not be a dominant strategy, because in practice it is not one:
risk accrues while an organization waits. The card therefore carries an interest
cost. Every threat left unmitigated at the end of a deferred round moves one
cell toward \emph{Very Likely} when the scenario is revisited, and the team must
state on the record which exposure it is accepting and for how long.

Drift alone is insufficient, because the matrix gives likelihood only three
states and a deferred threat therefore reaches \emph{Very Likely} within two
engagements with nowhere further to go. Capping it there would teach that
ignoring a risk plateaus, which is false. A threat at \emph{Very Likely} that
survives another engagement unmitigated instead \emph{realizes}: the facilitator
announces that it has occurred, and the outcome is determined not by chance but
by what the team purchased in earlier engagements. Whether the incident is
discovered in days or in months depends on whether logging or managed detection
was ever bought; whether data returns depends on backups; whether the
organization knows who to call depends on incident response planning; whether
the cost is survivable depends on insurance. Severity is read from the
\emph{Human Impact} card in play, and the team narrates the consequence to the
room in the client's terms. Appendix~\ref{app:realization} gives the resolution
table.

Realization exists to correct a structural bias in the deck. Twenty-one of its
twenty-eight cards are preventive, and novices spend on prevention first; a team
that never experiences an incident never learns why detection, response, and
recovery were worth points in the first place. This supplies that lesson without
adding a card, by making the earlier portfolio determine how bad the day is.

Run across a semester---the same client organization revisited in weeks 3, 8,
and 14---this yields a longitudinal artifact. Whether a team's week-14 portfolio
differs in quality from its week-3 portfolio is a learning measure that requires
no control group, and we intend to collect it alongside the study in
Section~\ref{sec:planned}.

\subsection{Optional Inject Expansion}
\label{sec:injects}

A recurring observation across our pilots is that nothing is at stake after
purchase. A team buys a portfolio, re-places the threat, and the round ends;
no choice is ever tested. The natural remedy is adversity, but adversity has to
be introduced carefully, because randomness that merely punishes teaches
nothing. A team that loses a card to a die roll has made no decision and has
nothing to reason about.

We therefore place adversity in a separate class of four \emph{inject} cards,
played by the facilitator rather than drawn by the team, and modeled on the
funding and staffing realities of the organizations clinic students serve
(Appendix~\ref{app:injects}). Each inject removes a capability for a stated
reason and requires the team to respond rather than simply absorb a loss: a
grant is not renewed and two points are withdrawn mid-round; a vendor
discontinues a product the team purchased; the one person who ran a control
leaves; or a restricted grant arrives that must be spent immediately on
something new. In each case the team must say what it gives up, what risk
returns, and what it would tell the client.

Two cautions govern their use. First, injects presuppose that teams have already
internalized the base loop, so we recommend them from the midpoint of a course
onward rather than in week one. Second, and more importantly for this paper,
\textbf{injects are excluded from the controlled study in
Section~\ref{sec:planned}}. Because they introduce variance that is not held
constant across teams, including them would confound the comparison between the
four-dimension and five-dimension conditions. They are a classroom expansion,
not part of the evaluated artifact.

\subsection{Why Collaborative Rather Than Competitive Play}

Teams of two to three play cooperatively against the scenario rather than
competitively against each other. This follows from the ICAP framework
\cite{chi2014icap}, which distinguishes interactive engagement---where learners
build on each other's contributions---from merely active engagement, and
associates the former with deeper outcomes. The mechanism we want is
articulation and challenge: a student who must say aloud why endpoint detection
beats segmentation for this threat, to a peer who may disagree, engages the
material differently from one who silently ranks cards.

The budget constraint supplies the productive friction that a purely
cooperative structure otherwise lacks. Pilot participants told us that too many
answers felt equally acceptable; six points and eight cards make that untrue.
Inter-team comparison at the end of Phase 3 supplies competitive motivation
without requiring a winner, since the interesting question is why two teams
facing the same threat bought different portfolios.

% =====================================================================
\section{Walkthrough}
\label{sec:walkthrough}

\subsection{Session A: Credential Phishing Against an Individual}

The drawn cards were \emph{Method: Manipulation or Coercion}, \emph{Impact:
Physical Wellbeing}, \emph{Resources: Money}, and \emph{Motivation: Desire or
Obsession}. The team constructed a scenario in which an attacker funds a
convincing replica of a social media login page, harvests a public figure's
credentials, and recovers a home address from the compromised account.

The team's discussion moved between digital compromise and physical safety, and
their ranking reflected that: Social Engineering Awareness first, on the
reasoning that recognition would have prevented the initial compromise;
Identity and Access Management second, to limit what stolen credentials reach;
Password and Authentication third, for the multi-factor requirement that would
have made the harvested password insufficient.

This session predates the version~2 deck and illustrates the coverage gap
directly. No Email and Browser Protections card existed, so DNS filtering
(safeguard 9.2)---a cost-1 \ig{1} control that plausibly blocks the replica
domain outright---was never on the table. The team reached for awareness
training partly because it was the only anti-phishing instrument available.
This is the clearest single piece of evidence we have that coverage analysis
changes what students consider.

\subsection{Session B: Abuse of a Publishing Workflow}

The drawn cards were \emph{Method: Processes}, \emph{Impact: Relationships},
\emph{Resources: Entitlement}, and \emph{Motivation: Politics}. The scenario
involved an attacker exploiting a debug-mode endpoint on a news site to publish
unauthorized articles attributed to a politician.

The team ranked Identity and Access Management first, Password and
Authentication second, and Penetration Testing third. They also raised audit
logging, defense in depth, deepfake detection, and endpoint detection, but set
them aside on the grounds that these address detection and response rather than
preventing initial access---an unprompted distinction between control functions
that we regard as the strongest qualitative signal from the pilots.

Under version~2 rules this session changes materially. Penetration Testing now
costs three of six points, forcing the team to state what it would give up;
Secure Configuration (safeguard 4.1, cost~1) becomes available and is arguably
the correct answer, since a debug endpoint in production is a configuration
failure; and Phase~3 would require the team to say what residual risk remains
after their purchases.

% =====================================================================
\section{Formative Evaluation}
\label{sec:eval}

\subsection{Purpose and Claims}

This evaluation is formative. Its purpose was to determine whether novices can
use the deck productively without facilitator expertise, how long play takes,
and whether the vocabulary is accessible. We do not claim learning gains, and
the design---a single session with a post-play survey and no control
condition---cannot support such a claim.
\todo{Insert IRB protocol number or exemption determination. Anonymize for
initial submission.}

\subsection{Participants}

Nineteen participants were recruited from a summer research program at a
comprehensive undergraduate institution. Table~\ref{tab:participants} reports
the composition.
\todo{Complete this table from your records. Reviewers of the earlier version
specifically requested discipline distribution, academic level, and prior
gaming experience; leaving cells blank is worse than reporting small counts.}

\begin{table}[h]
\centering
\caption{Formative evaluation participants ($N=19$).}
\label{tab:participants}
\small
\begin{tabular}{@{}l r@{}}
\toprule
\textbf{Characteristic} & \textbf{$n$} \\
\midrule
\multicolumn{2}{@{}l}{\textit{Discipline}} \\
\quad Computer science & \textcolor{red}{--} \\
\quad Psychology & \textcolor{red}{--} \\
\quad Biology & \textcolor{red}{--} \\
\quad Physics & \textcolor{red}{--} \\
\quad Mathematical sciences & \textcolor{red}{--} \\
\quad Other & \textcolor{red}{--} \\
\midrule
\multicolumn{2}{@{}l}{\textit{Academic level}} \\
\quad Undergraduate & \textcolor{red}{--} \\
\quad Graduate & \textcolor{red}{--} \\
\quad Faculty & \textcolor{red}{--} \\
\midrule
\multicolumn{2}{@{}l}{\textit{Prior security coursework}} \\
\quad None & \textcolor{red}{--} \\
\quad One course & \textcolor{red}{--} \\
\quad More than one & \textcolor{red}{--} \\
\midrule
\multicolumn{2}{@{}l}{\textit{Self-rated security knowledge (1--5)}} \\
\quad Mean (SD) & 2.53 (\textcolor{red}{--}) \\
\bottomrule
\end{tabular}
\end{table}

\subsection{Procedure}

Six teams of two to three played simultaneously in a single room, each
facilitated by a member of the research team. Teams received a hypothetical
organization and an attack scenario, then used the deck to select and rank the
three most important defensive measures. Sessions ran seven to ten minutes
excluding debrief and survey. Play was collaborative; no competitive scoring
was used, and the budget mechanic described in Section~\ref{sec:gameplay} was
not yet implemented.

\subsection{Results}

Seventeen of nineteen participants agreed or strongly agreed that the cards
helped them learn. Ease of use was the highest-rated dimension at a mean of
4.42 on a five-point scale. Perceived sufficiency of card content scored 3.95.
Figure~\ref{fig:likert} reports the full distributions.
\todo{Replace the four pie charts from the earlier draft with a single diverging
stacked bar chart, and report counts rather than percentages. With $N=19$,
``31.6\%'' asserts a precision the sample does not have; ``6 of 19'' is both
honest and easier to read.}

\begin{figure}[h]
\centering
\figplaceholder{3.4cm}{Figure placeholder.\\[0.4em]
Insert a diverging stacked bar chart of the four survey items,
with response counts labeled on each segment.\\[0.4em]
Replace this box with an includegraphics call
pointing at likert.pdf, sized to the column width.}
\caption{Post-play survey responses ($N=19$) for the four evaluation items,
reported as counts on a five-point agreement scale.}
\label{fig:likert}
\end{figure}

Open-ended responses converged on three themes. Participants valued the
definitions and examples on the card backs and the non-specialist vocabulary,
which speaks directly to goal G1. Participants wanted feedback on whether their
rankings were correct, and noted that several solutions seemed equally valid.
Participants found fifteen cards to be many, and some card descriptions
overlapping.

The second and third themes are the origin of the version~2 design. The request
for feedback and the complaint that too many answers were acceptable are the
same observation: unconstrained ranking has no wrong answers, so it generates
agreement rather than argument. The budget mechanic and the residual-risk phase
were introduced specifically to make choices costly and therefore contestable.
The deck-size concern motivated the tiering and the eight-card hand described in
Section~\ref{sec:design}.

\subsection{What This Evaluation Establishes}

Three things, modestly stated. The activity is feasible in under ten minutes,
which matters for crowded syllabi. It is usable by participants with a mean
self-rated security knowledge of 2.53 out of 5 and no facilitator expertise
beyond reading the scenario. And it generates the design signal that produced
version~2. It establishes nothing about learning.

% =====================================================================
\section{Planned Classroom Evaluation, September 2026}
\label{sec:planned}

We report the design of the controlled study in advance of running it, both so
that it can be criticized before rather than after data collection and so that
the claims in this paper remain clearly separated from the claims that study
would license.

\textbf{Setting and participants.} An undergraduate cybersecurity clinic course,
$N \approx$ \todo{expected enrollment}, during the first three weeks of the
Fall 2026 term.
\todo{IRB submission must precede the term; note approval status here.}

\textbf{Design.} A counterbalanced within-subjects crossover. Each team plays
two scenarios of matched difficulty. In one, teams use the four-dimension deck
and are asked to discuss defenses without Solutions cards; in the other, they
use the five-dimension deck with the budgeted V6 loop. Order and
scenario--condition pairing are counterbalanced across teams, so each team
serves as its own control and scenario effects are separated from condition
effects.

\textbf{Measures.}
\begin{description}[leftmargin=1em,style=nextline]
\item[M1. Control-selection appropriateness.] Purchased portfolios scored
against an expert key derived from CIS safeguard--threat mappings. Two raters,
blind to condition; agreement reported as Cohen's $\kappa$
\cite{landis1977kappa}.
\item[M2. Risk-matrix delta.] Magnitude and direction of movement between
pre- and post-control placements, and whether the team's stated residual risk
is consistent with the controls purchased.
\item[M3. Justification quality.] Each stated rationale scored 0--3 on a rubric
ranging from assertion without reasoning to a threat-specific mechanism
correctly linked to a safeguard.
\item[M4. Mitigation breadth.] Count of distinct defensive measures raised per
scenario, the direct test of whether Solutions cards expand or merely redirect
discussion.
\item[M5. Self-efficacy.] Pre/post administration of a security self-efficacy
scale \cite{bandura1997selfefficacy}.
\todo{Select and cite a validated instrument rather than authoring items.}
\item[M6. Cognitive load.] Short-form NASA-TLX \cite{hart1988nasatlx} to test
whether the expanded deck and budget mechanic overload novices, given the
pilot signal on deck size.
\item[M7. Card comprehension.] A brief matching task, card face to correct
definition, testing the icon revision that we currently support only by
facilitator observation.
\end{description}

\textbf{Analysis.} Paired comparisons on M1--M4 using Wilcoxon signed-rank tests
given the expected sample size and ordinal measures; descriptive reporting for
M5--M7. We will report effect sizes and confidence intervals rather than
significance alone, and will treat the study as adequately powered only for
large effects.

\textbf{Anticipated threats to validity.} Team size varies with attendance.
Facilitators are not blind to condition. Scenario difficulty is matched by
expert judgment rather than empirically. Novelty effects favor whichever
condition is played second. We will report all four rather than control for
them, since a single-course study cannot eliminate them.

% =====================================================================
\section{Discussion and Limitations}
\label{sec:discussion}

\textbf{What the artifact is for.} The deck is an on-ramp, not a simulation. It
does not carry the technical depth required to conduct a real risk assessment,
and a student who plays it well is not thereby competent to advise a client. Its
function is to give a novice a vocabulary, a framework identifier they can look
up, and forty minutes of practice at the specific cognitive move---from stated
threat to proportionate, affordable control---that the rest of a clinic
curriculum builds on.

\textbf{The framework is a commitment, not a decoration.} Anchoring to CIS
Controls buys traceability and an externally justified cost model, and it costs
generality. A course oriented to the NIST Cybersecurity Framework
\cite{nist_csf2} or working toward CSEC2017 knowledge areas \cite{csec2017}
would want a different mapping on the card backs. We regard the mapping
procedure in Section~\ref{sec:gaps} as portable even where the specific control
identifiers are not.

\textbf{Why there is no sixth dimension.} A natural proposal is to add an
Impact dimension so that students see consequences alongside threats,
mitigations, and cost. We considered and rejected it, because the original deck
already contains one: \emph{Human Impact} is one of the four dimensions, and a
fifth impact category would duplicate a taxonomy that is already good. What was
genuinely missing was not a category but a moment---the point at which
consequence is determined by the team's own choices rather than described in the
abstract. Realization supplies that moment, and binding severity placement to
the Human Impact card connects the existing dimension to the matrix. Both are
rules rather than cards, which keeps the deck from growing to solve a problem
that was never about card count.

\textbf{Limitations.} The formative evaluation had no control condition, no
pre/post measure, and a single short session, and it evaluated version~1 of the
deck rather than the version described here. Participants were volunteers in a
research program, which likely selects for engagement. The icon revision is
supported by facilitator observation only. The cost model, while derived from
Implementation Groups, still involves judgment where a card spans tiers, and
reasonable designers would tier Penetration Testing or Network Segmentation
differently. Finally, the version~2 deck has not yet been played in a
classroom; Section~\ref{sec:planned} is a plan, not a result.

\textbf{Future work.} Three directions seem worthwhile beyond the September study. A sixth dimension representing
the \emph{defender's} resources would let scarcity vary by organization rather
than by rule. A scenario deck of named incident types would standardize scenarios
across sessions and remove a confound from any multi-team study. And
\emph{standing risks}---two chronic problems the client already has, placed on
the matrix alongside the presenting incident and competing for the same
budget---would test something the current single-threat loop cannot: whether
students recognize that a control touching several problems is worth more than
one addressing only today's. That is the argument for foundational controls, and
at present nothing in the design rewards it. All three are untested.

% =====================================================================
\section{Conclusion}

The Security Cards are an effective instrument for the task they were designed
for and stop deliberately at the boundary of that task. In cybersecurity clinic
education, where students must convert threat understanding into affordable
recommendations for organizations that cannot afford much, that boundary falls
in the wrong place. We have presented a Solutions dimension of twenty-eight
cards---twenty-one defensive controls mapped to CIS Controls v8.1 at safeguard
level and seven zero-cost principle modifiers, a cost model derived from
Implementation Group tier rather than designer judgment, a coverage analysis
that identified and closed five foundational gaps in our own earlier deck, and
a gameplay framework whose core loop scores residual risk before and after
control selection. A formative evaluation with 19 participants established
feasibility and accessibility; a controlled classroom study is scheduled for
September 2026 and its design is reported here in full. The deck is released
openly, and we would welcome its use, criticism, and re-tiering by other clinic
programs.

% =====================================================================
\section*{Generative AI Disclosure}
\todo{CISSE requires disclosure of any use of generative AI or automated tools
in creating this content. State here what was used and for what. AI tools may
not be listed as authors.}

\section*{Acknowledgments}
\todo{Omit from the anonymized submission; restore for camera-ready.}

% =====================================================================
% APPENDIX: full card copy
% =====================================================================

\input{cards.tex}

% =====================================================================
% BIBLIOGRAPHY
%
% Default: a self-contained thebibliography, so the document compiles
% with a single pdflatex pass and no bibtex run.
%
% To use BibTeX instead (recommended on Overleaf, where references.bib
% is already provided), comment out the \input line and uncomment the
% two lines beneath it.
% =====================================================================

\input{references.tex}

% \bibliographystyle{plain}   % or IEEEtran if available
% \bibliography{references}

\end{document}
